% =============================================================================
% Project : THE_PHILOSOPHY_OF_VIABILITY
% Book    : Book I — Ontology of Viability
% Book ID : VS.B01
%
% File    : 4_ASSETS/diagrams/
%           VS.B01.P01.C02.S08.SS05.P004.DG01.continuity_and_invariance_relations.tex
%
% Asset   : DG01 — Continuity and Invariance Relations
% Location: Chapter 02 / Section 08 / Subsection 05 / Paragraph 004
%
% Purpose :
%   To distinguish the analytical questions of invariance
%   and continuity and to display their possible coexistence
%   or independent occurrence.
%
% Requires (preamble):
%   \usepackage{tikz}
%   \usetikzlibrary{arrows.meta, calc, positioning}
% =============================================================================

\begin{figure}[htbp]
    \centering

    \begin{tikzpicture}[
        font=\small,
        node distance=1.6cm and 2cm,
        % --- base box ---
        box/.style={
            draw, rounded corners=3pt, align=center,
            text width=4.8cm, minimum height=1.6cm,
            inner sep=6pt, line width=0.6pt
        },
        % --- node roles ---
        concept/.style={box, fill=blue!8,  draw=blue!55!black},
        outcome/.style={box, fill=gray!10, draw=black!60},
        joint/.style  ={box, fill=green!8, draw=green!40!black, text width=5cm},
        % --- arrows ---
        relation/.style={-{Stealth[length=2.5mm]}, thick},
        optional/.style={-{Stealth[length=2.5mm]}, thick, dashed},
        % --- arrow labels ---
        elabel/.style={font=\scriptsize\itshape, anchor=south, inner sep=2pt}
    ]

    % ---------------------------------------------------------------
    % Row 1: analytical questions
    % ---------------------------------------------------------------
    \node[concept] (invariance) {
        \textbf{Invariance}\\[3pt]
        Preservation of a specified determination under a specified transformation
    };
    \node[concept, right=of invariance] (continuity) {
        \textbf{Continuity}\\[3pt]
        A determinate relation connecting successive states or phases
    };

    % ---------------------------------------------------------------
    % Row 2: outcomes
    % ---------------------------------------------------------------
    \node[outcome, below=of invariance] (preserved) {
        \textbf{Preservation established}\\[3pt]
        A relevant determination remains preserved
    };
    \node[outcome, below=of continuity] (connected) {
        \textbf{Connection established}\\[3pt]
        A relevant relation between states or phases is established
    };

    % ---------------------------------------------------------------
    % Row 3: possible coexistence (centered between the two outcomes)
    % ---------------------------------------------------------------
    \node[joint, anchor=north] (joint) at
        ($(preserved.south)!0.5!(connected.south)+(0,-2.2cm)$) {
        \textbf{Possible coexistence}\\[3pt]
        Both may apply, but neither proves the other
    };

    % ---------------------------------------------------------------
    % Note
    % ---------------------------------------------------------------
    \node[below=0.9cm of joint, text width=11cm, align=center,
          font=\footnotesize\itshape] (limit) {
        Neither invariance nor continuity independently establishes
        the numerical identity of an existent.
    };
    \draw[black!40, line width=0.4pt]
        ($(limit.north west)+(1.5,0.3)$) -- ($(limit.north east)+(-1.5,0.3)$);

    % ---------------------------------------------------------------
    % Arrows
    % ---------------------------------------------------------------
    \draw[relation] (invariance.south) -- (preserved.north);
    \draw[relation] (continuity.south) -- (connected.north);

    % down, across, down into the top of the joint box
    \draw[optional] (preserved.south) -- ++(0,-0.9)
        -| ([xshift=1.1cm]joint.north west);
    \draw[optional] (connected.south) -- ++(0,-0.9)
        -| ([xshift=-1.1cm]joint.north east);

    % labels sit above the horizontal segments (no background needed)
    \node[elabel] at ($(preserved.south)+(1.8,-0.9)$) {may be relevant to};
    \node[elabel] at ($(connected.south)+(-1.8,-0.9)$) {may be examined with};

    \end{tikzpicture}

    \caption{Invariance and continuity address distinct analytical
    questions. Their coexistence is possible, but neither relation
    independently establishes the other or numerical identity.}
    \label{fig:continuity-and-invariance-relations}
\end{figure}

% =============================================================================
% END OF DIAGRAM DG01
% =============================================================================